% =========================================================================
%  src/goto-btor2 形式化语义与可靠性定理
%
%  本文档形式化描述 goto-btor2 转换器将 goto 程序翻译为 BTOR2 有限状态
%  迁移系统的语义，给出编码所满足的可靠性定理，并讨论语义层次、内存
%  模型支持程度。
%
%  编译:  pdflatex goto-btor2-semantics.tex  (需两遍)
% =========================================================================
\documentclass[11pt,a4paper]{article}

\usepackage[UTF8]{ctex}
\usepackage[margin=2.5cm]{geometry}
\usepackage{amsmath,amssymb,amsthm}
\usepackage{stmaryrd}
\usepackage{mathtools}
\usepackage{mathpartir}
\usepackage{enumitem}
\usepackage{booktabs}
\usepackage{array}
\usepackage{hyperref}
\usepackage{xcolor}
\usepackage{framed}

\hypersetup{colorlinks=true,linkcolor=blue!70!black,urlcolor=blue!60!black}
\definecolor{shadecolor}{rgb}{0.95,0.95,0.92}

\newtheorem{definition}{定义}[section]
\newtheorem{lemma}{引理}[section]
\newtheorem{theorem}{定理}[section]
\newtheorem{corollary}{推论}[section]
\newtheorem{remark}{注记}[section]

\newcommand{\Loc}{\mathit{Loc}}
\newcommand{\State}{\mathit{State}}
\newcommand{\Var}{\mathit{Var}}
\newcommand{\Val}{\mathit{Val}}
\newcommand{\Input}{\mathit{Input}}
\newcommand{\Reach}{\mathit{Reach}}
\newcommand{\Init}{\mathit{Init}}
\newcommand{\Trans}{\mathit{Trans}}
\newcommand{\Bad}{\mathit{Bad}}
\newcommand{\Cons}{\mathit{Constraint}}
\newcommand{\eval}[1]{\llbracket #1 \rrbracket}
\newcommand{\guard}{\chi}
\newcommand{\nextv}[1]{\mathit{next}(#1)}
\newcommand{\nextpc}{\nextv{\mathit{pc}}}
\newcommand{\depth}{\mathit{depth}}
\newcommand{\True}{\mathit{true}}
\newcommand{\False}{\mathit{false}}
\newcommand{\ite}{\mathop{\mathit{ite}}}

\title{\bfseries src/goto-btor2 形式化语义与可靠性定理}
\author{CBMC 项目文档}
\date{\today}

\begin{document}
\maketitle
\tableofcontents

% =========================================================================
\section{引言}
% =========================================================================

本文档对 \texttt{src/goto-btor2} 模块进行形式化描述。该模块将 CBMC 的
goto 中间表示（一种简化的命令式程序 IR）翻译为 BTOR2 格式的有限状态
迁移系统（Finite Transition System, FTS），供硬件模型检测器
（如 btormc、Pono）进行安全性质检查。

本文档回答以下问题：
\begin{enumerate}[label=(\arabic*)]
  \item 转换器生成的 BTOR2 模型的精确形式语义是什么？（\S\ref{sec:semantics}）
  \item 该翻译是否可靠（sound）？在什么条件下成立？（\S\ref{sec:soundness}）
  \item 编码属于哪种语义层次？能否导出霍尔推理规则？（\S\ref{sec:semantic-level}）
  \item 内存模型的支持程度如何？（\S\ref{sec:memory-model}）
\end{enumerate}

% =========================================================================
\section{形式化语义}
\label{sec:semantics}
% =========================================================================

\subsection{标量数据}
\label{sec:scalar}

\paragraph{类型与分析。}
C 的标量类型包括布尔（\texttt{\_Bool}）、字符（\texttt{char}）、
整数（\texttt{short}、\texttt{int}、\texttt{long}、\texttt{long long}，
及其 \texttt{unsigned} 变体）、枚举（\texttt{enum}）和浮点
（\texttt{float}、\texttt{double}、\texttt{long double}）。
这些类型在 C 语义中有两个核心属性：\emph{位宽}（决定值域大小）
和\emph{符号性}（决定算术、比较、扩展的行为）。
此外，C 标准将大量边界情况——有符号整数溢出、除以零、过大移位量——
留给\emph{实现定义}或\emph{未定义行为}，
这意味着同一段 C 代码在不同编译器上可能产生不同结果。

\paragraph{BTOR2 的能力。}
BTOR2 只有一个标量类型：位向量 $\mathrm{bv}\langle n \rangle$。
它没有 signed/unsigned 之分——类型标签在 sort 层面完全丢失。
但 BTOR2 对同一个位向量同时提供有符号和无符号的操作原语
（如 $\mathit{sdiv}$ / $\mathit{udiv}$、$\mathit{slt}$ / $\mathit{ult}$），
因此符号性不在 sort 里记录，而在\emph{操作指令}层面区分。
位向量的算术语义是确定的：加法、乘法按 $2^n$ 回绕，
不存在 C 的未定义行为——溢出就是回绕，结果总是确定的。

\paragraph{困难。}
核心困难不在于"位向量能不能表示整数"——显然可以——
而在于\emph{如何在每一个操作上恢复正确的 C 语义}。
具体来说：

（1）\textbf{符号性的恢复}。
同一个 $\mathrm{bv}\langle 32 \rangle$ 上的除法，
\texttt{int} / \texttt{int} 应该用 $\mathit{sdiv}$，
\texttt{unsigned} / \texttt{unsigned} 应该用 $\mathit{udiv}$。
编码器必须在每个操作处查询操作数类型的符号性，
选择正确的 BTOR2 原语。

（2）\textbf{类型转换的语义}。
窄类型转宽类型时，\texttt{signed char} $\to$ \texttt{int}
需要符号扩展（$\mathit{sext}$，复制符号位），
\texttt{unsigned char} $\to$ \texttt{int} 需要零扩展（$\mathit{uext}$，高位补零）。
宽类型转窄类型时，C 标准规定截取低位，这对有符号和无符号是一致的。
但同宽度的转换（如 \texttt{unsigned} $\to$ \texttt{int}）是位模式的重新解释，
不做任何变换。

（3）\textbf{未定义行为的处理}。
除以零在 C 中是未定义行为；在 BTOR2 中 $\mathit{udiv}$ 以零为除数
返回全 1（$\mathit{ones}$），是一个确定的（但非 C 的）值。
过大移位量（$\geq$ 位宽）在 C 中未定义，
BTOR2 的移位是全函数（结果为 0）。
这些差异如果不处理，可能导致编码器遗漏真实的错误。

（4）\textbf{浮点的精确性}。
C 浮点遵循 IEEE 754 标准，包含 NaN、$\pm\infty$、次正规数、
四种舍入模式。
BTOR2 没有浮点 sort，只有一个不透明的位向量。
要把 \texttt{float} 的加法翻译到 BTOR2，
就需要把整个 IEEE 754 加法器的电路展开成布尔逻辑。

\paragraph{方法。}
我们为标量类型定义类型翻译函数：
\begin{align*}
\lceil \texttt{\_Bool} \rceil &= \mathrm{bv}\langle 1 \rangle \\
\lceil \texttt{char} \rceil &= \mathrm{bv}\langle 8 \rangle, \quad
\lceil \texttt{short} \rceil = \mathrm{bv}\langle 16 \rangle \\
\lceil \texttt{int} \rceil &= \mathrm{bv}\langle 32 \rangle, \quad
\lceil \texttt{long long} \rceil = \mathrm{bv}\langle 64 \rangle \\
\lceil \texttt{unsigned}\;\tau \rceil &= \lceil \tau \rceil
  \quad \textit{(位宽相同)} \\
\lceil \texttt{enum}\;E \rceil &= \mathrm{bv}\langle w_E \rangle
  \quad \textit{($w_E$ = 底层整数类型的位宽)} \\
\lceil \texttt{float} \rceil &= \mathrm{bv}\langle 32 \rangle, \quad
\lceil \texttt{double} \rceil = \mathrm{bv}\langle 64 \rangle
\end{align*}
位宽来自 CBMC 的目标平台配置（如 LP64 上 \texttt{long} = 64 位，
ILP32 上为 32 位），编码器直接读取，不做二次推导。
每个标量变量 $v : \tau$ 产生一个 BTOR2 状态
$x_v : \mathit{state}\;\lceil \tau \rceil$。

\textbf{符号性在操作层面恢复。}
编码器维护一个判定函数 $\mathit{is\_signed}(\tau)$，
在遇到除法、取模、比较、宽化转换时查询：
\begin{itemize}
\item 除法 / 取模：$\mathit{is\_signed}$ 为真用 $\mathit{sdiv}$ / $\mathit{srem}$，
  否则用 $\mathit{udiv}$ / $\mathit{urem}$。
  这里用的是 $\mathit{srem}$（截断取模，符号跟随被除数）而非 $\mathit{smod}$（向下取整），
  与 C99 的 \texttt{\%} 语义一致。
\item 比较：有符号用 $\mathit{slt}$ / $\mathit{slte}$ / $\mathit{sgt}$ / $\mathit{sgte}$，
  无符号用 $\mathit{ult}$ / $\mathit{ulte}$ / $\mathit{ugt}$ / $\mathit{ugte}$。
  等于 / 不等于不区分符号性（$\mathit{eq}$ / $\mathit{neq}$）。
\item 宽化转换：$\mathit{is\_signed}$ 为真用 $\mathit{sext}$，否则 $\mathit{uext}$。
  窄化转换统一用 $\mathit{slice}$ 取低位。
\end{itemize}
符号性的判定依据操作数的 C 类型 ID
（CBMC 的 \texttt{ID\_signedbv} / \texttt{ID\_unsignedbv}），
而非任何 BTOR2 层面的信息。

\textbf{未定义行为的处理。}
CBMC 前端在每个除法前自动注入 \texttt{assert(divisor != 0)} 断言，
在每个算术运算前注入溢出检查断言。
这些断言被编码为独立的 BTOR2 bad 属性。
因此除以零和算术溢出在\emph{断言层}被检测，
而非在操作层处理——如果除数为零，
对应的 bad 属性会先被触发，
模型检查器在评估商之前就报告了违规。

\textbf{浮点的位级展开。}
浮点值存储为 IEEE 754 打包位向量。
所有浮点运算通过适配器 \texttt{btor2\_propt} 复用 CBMC 的
\texttt{float\_utilst}——CBMC 的参考 IEEE 754 实现——
将加减乘除、比较、\texttt{isnan}、\texttt{isinf} 等操作
展开为等价的布尔逻辑电路。
具体来说，\texttt{float\_utilst} 将浮点值拆成符号位、指数、尾数
三个位向量字段，用整数操作模拟 IEEE 754 的对齐、舍入、规格化过程。
适配器为每个新引入的布尔变量创建一个 1-bit BTOR2 input，
为每个等式约束创建一个 BTOR2 constraint 节点。
这些 constraint 是\emph{全局}的（不受 PC 守卫限制），
因为它们是位展开电路的重言式，不是程序假设。

\paragraph{效果与保真度。}
对整数和布尔类型，在 C 的"有符号溢出为回绕"这一实现定义行为下，
编码是\emph{精确的}：BTOR2 表达式与 C 表达式在每一位上的值一致。
CBMC 注入的溢出检查断言在 BTOR2 中成为独立 bad 属性，
因此如果目标程序依赖未定义行为（如 signed overflow），
模型检查器会报告溢出断言失败，而不是产生错误的计算结果。

浮点编码在 IEEE 754 标准语义下也是精确的：
NaN 传播、$\pm\infty$、次正规数、四种舍入模式均正确处理。

\paragraph{局限性。}
浮点转整数仅支持向零舍入（round-to-zero），
其他舍入模式下的 \texttt{float} $\to$ \texttt{int} 转换不受支持，
回退为无约束近似。
\texttt{long double} 的位宽和格式取决于目标平台
（80 位 x87 扩展精度、128 位四精度等），
编码器按配置展开，但不保证所有平台组合都已测试。


% =====================================================================
\subsection{复合数据：结构体}
\label{sec:struct}

\paragraph{类型与分析。}
C 的 \texttt{struct} 是一组命名成员的聚合，
成员可以是标量、嵌套结构体、数组、指针，甚至是灵活数组成员（C99）。
结构体的关键语义属性包括：
成员的\emph{布局}（偏移量、对齐、padding）、
\emph{整体赋值}（$s_1 = s_2$ 复制所有成员）、
\emph{按成员访问}（$s.field$）、
以及\emph{数组的关系}（结构体数组的每个字段构成一个"列存"视图）。

\paragraph{BTOR2 的能力。}
BTOR2 没有 struct sort。
它只有位向量和数组，且数组的元素只能是位向量或（嵌套的）数组。
没有任何机制可以声明一个包含异构成员的复合类型。

\paragraph{困难。}
（1）\textbf{异构聚合的表示}。
一个 \texttt{struct \{int x; float y; char z;\}} 有三个类型不同的成员，
BTOR2 无法用一个状态变量表示它们。

（2）\textbf{整体赋值}。
\texttt{s1 = s2} 需要同时复制所有成员，
如果每个成员是独立状态，就需要生成多个并行的 next 更新。

（3）\textbf{结构体数组}。
\texttt{struct S arr[N]} 既要支持 \texttt{arr[i].x}（访问第 $i$ 个元素的 $x$ 字段），
又要支持 \texttt{arr[i] = some\_struct}（整体赋值）。
如果简单地把每个 $(i, \mathit{field})$ 组合拆成独立标量，
动态索引 $a[i].x$ 需要一个 $O(N)$ 的 ITE 链。

\paragraph{方法：字段级展平。}
我们将结构体递归拆解为独立的状态分量。
定义展平投影 $\pi$：
\[
\pi(\texttt{struct}\;\{f_1{:}\tau_1, \ldots, f_n{:}\tau_n\})
= \{(f_1, \pi(\tau_1)), \ldots, (f_n, \pi(\tau_n))\}
\]
其中 $\pi(\tau)$ 对标量类型返回自身，
对嵌套结构体递归展开，对数组另行处理（见下文）。
padding 成员（CBMC 标记为 \texttt{ID\_C\_is\_padding}）在投影时丢弃。

变量 $s : \texttt{struct}\;\{x{:}\texttt{int},\;y{:}\texttt{int}\}$
产生两个独立 BTOR2 状态：
$s.x : \mathit{state}\;\mathrm{bv}\langle 32 \rangle$，
$s.y : \mathit{state}\;\mathrm{bv}\langle 32 \rangle$。

\textbf{成员访问} $s.f$ 翻译为对展平后状态 $x_{s.f}$ 的直接读写，
不需要任何偏移计算或内存布局。

\textbf{整体赋值} $s_1 = s_2$ 被递归投影为逐字段赋值：
$s_1.f_i \leftarrow s_2.f_i$，每个字段独立产生一个 next 更新。
若右值是结构体常量 $\{v_1, \ldots, v_n\}$，
则按位置拆分：$s_1.f_i \leftarrow v_i$。

\textbf{结构体数组} $\texttt{struct S arr[}N\texttt{]}$
按\emph{字段投影}为独立数组：
$\mathit{arr.x}$ 是一个元素类型为 \texttt{int} 的数组，
$\mathit{arr.y}$ 是另一个。
访问 $\mathit{arr[}i\mathit{]}.x$ 翻译为对 $\mathit{arr.x}$ 的索引读。
整体赋值 $\mathit{arr[}i\mathit{]} = \mathit{val}$
被投影为逐字段写入：
$\mathit{arr.x[}i\mathit{]} \leftarrow \mathit{val.x}$，
$\mathit{arr.y[}i\mathit{]} \leftarrow \mathit{val.y}$。

\paragraph{效果。}
结构体的字段访问和整体赋值都是\emph{精确的}。
展平后的每个状态变量与 C 语义中对应字段的值完全一致。
padding 的丢弃不影响语义，因为 C 标准规定 padding 的值不可观察。

\paragraph{局限性。}
展平破坏了结构体的\emph{布局信息}：
无法通过指针做字节级访问（如 \texttt{memcpy} 结构体、
通过 \texttt{char*} 逐字节读取）。
但这属于指针层的限制（见第~\ref{sec:pointer}~节），
结构体字段级操作本身没有精度损失。


% =====================================================================
\subsection{复合数据：联合体}
\label{sec:union}

\paragraph{类型与分析。}
C 的 \texttt{union} 让多个成员\emph{共享同一块存储}。
关键语义规则是\emph{活跃成员}（active member）：
写入 $u.i$ 使 $i$ 成为活跃成员，
随后读取 $u.f$（另一个成员）在 C 标准中是未定义行为
（但大多数实现支持 type-punning，即按位模式重新解释）。
联合体的大小等于最大成员的大小。

\paragraph{BTOR2 的能力与困难。}
BTOR2 无法表达"同一存储的多重视图"。
如果把联合体的每个成员建为独立状态变量，就丢失了共享存储的语义——
写入一个成员不会改变另一个成员的值，
而 C 语义中它们应该共享底层位模式。

\paragraph{方法：逐成员独立状态。}
当前实现将联合体的每个成员投影为独立的状态变量，
与结构体的投影逻辑\emph{完全相同}：
\[
\pi(\texttt{union}\;\{m_1{:}\tau_1, \ldots, m_k{:}\tau_k\})
= \{(m_1, \pi(\tau_1)), \ldots, (m_k, \pi(\tau_k))\}
\]
即 \texttt{union \{int i; float f;\} u} 产生两个独立状态
$u.i$ 和 $u.f$，它们之间没有任何约束。

\paragraph{效果。}
这是一个\emph{有局限的近似}。
对只通过活跃成员访问的联合体（写 $u.i$，读 $u.i$），
编码是精确的——因为只涉及一个状态变量。
但对 type-punning 场景（写 $u.i$，读 $u.f$），
两个状态变量互不影响，读到的 $u.f$ 值与写入的 $u.i$ 无关，
是错误的。

\paragraph{局限性。}
联合体的 active-member 语义\emph{未保留}。
不存在判别器（discriminator）状态，
不存在成员间的共享存储约束。
论文中应将联合体标注为"有限支持"，
并说明 type-punning 场景下的不可靠性。
位宽计算取第一个成员的位宽（\texttt{get\_width} 的实现），
这是实现上的约定，不是 C 标准的要求。


% =====================================================================
\subsection{复合数据：数组}
\label{sec:array}

\paragraph{类型与分析。}
C 数组是相同类型元素的连续序列。
关键操作包括：
按\emph{动态索引}读写（$a[i]$，$i$ 是运行时变量）、
\emph{整体赋值}（不支持，但初始化可以）、
\emph{多维数组}（$a[i][j]$，行优先布局）、
\emph{部分初始化}（\texttt{int a[3] = \{1\}}，剩余元素为零）。
C 数组不进行边界检查——越界访问是未定义行为。

\paragraph{BTOR2 的能力。}
BTOR2 有原生的数组 sort $\mathit{array}\langle n, m \rangle$
（索引为 $n$ 位位向量，元素为 $m$ 位位向量），
以及 $\mathit{read}$ 和 $\mathit{write}$ 操作。
BTOR2 数组是\emph{全函数}：对任何索引（包括越界索引）都返回一个值。
BTOR2 数组状态变量\emph{不支持 init}——初始内容无约束。
数组 sort 只支持一维，元素不能是另一个数组 sort（但可以是位向量）。

\paragraph{困难。}
（1）\textbf{初始值}。
BTOR2 数组状态没有 init，初始内容是任意值。
但 C 要求未显式初始化的静态数组为零，
部分初始化时剩余元素为零。

（2）\textbf{多维数组}。
BTOR2 只有一维数组，
C 的 \texttt{int a[N][M]} 需要线性化为一维。

（3）\textbf{性能}。
BTOR2 的 array sort 对 SAT/SMT 求解器来说比较昂贵。
IC3 等 IC3-style 归纳推理引擎在处理大数组时性能急剧下降。
如果一个小数组只有 4 个元素，
用 4 个独立标量状态 + ITE 选择器远比一个 array sort 高效。

（4）\textbf{越界}。
BTOR2 的 read/write 是全函数，
不会因越界访问而"报错"，
因此无法直接检测数组越界。

\paragraph{方法：两条路径。}

\textbf{路径 A — 小数组标量化}（元素数 $N \leq 256$ 且元素类型受支持）：
数组被拆成 $N$ 个独立标量状态
$x_{a\$0}, \ldots, x_{a\$(N{-}1)}$，
每个状态有独立的 init（零初始化或静态初值）。
动态索引读 $a[i]$ 翻译为嵌套条件选择：
\[
\mathcal{E}\llbracket a[i] \rrbracket
= \ite(i{=}0,\; x_{a\$0},\; \ite(i{=}1,\; x_{a\$1},\; \cdots))
\]
写入 $a[i] := e$ 为每个 $x_{a\$j}$ 生成条件更新：
$\mathit{next}(x_{a\$j}) = \ite(\mathit{at}(\ell) \wedge i{=}j,\; e,\; \cdots)$。
构建复杂度为 $O(N)$。

\textbf{路径 B — BTOR2 array state}（大数组或元素类型需要 array sort）：
\[
\lceil \texttt{T}[N] \rceil_{\mathrm{array}}
= \mathit{state}\;\mathit{array}\langle k,\; |\lceil T \rceil| \rangle
\]
索引位宽 $k$ 取 C 索引类型的位宽（通常 32 或 64），
若不可用则取 $k = \lceil \log_2 N \rceil$。
多维数组先线性化：
$\texttt{T}[N][M] \to \texttt{T}[N{\times}M]$，
$a[i][j] \to a[i \times M + j]$。
读写使用 BTOR2 原生操作：
$\mathcal{E}\llbracket a[i] \rrbracket = \mathit{read}(\mathit{state}[a],\; i)$，
$\mathit{next}(\mathit{state}[a]) = \mathit{write}(\mathit{state}[a],\; i,\; e)$。

\textbf{初始化处理。}
标量化数组：每个元素状态走标量的初始化逻辑（零或静态初值）。
部分初始化 \texttt{int a[3] = \{1\}} 的处理分两阶段：
先对全部元素写入零，再按操作数列表写入显式初值，
与 C 的聚合体零初始化语义一致。
BTOR2 array state：不生成 init，初始内容无约束——
这是一个上近似，意味着如果程序依赖数组初始值为零
（如全局数组），模型允许任意初值，可能导致虚假反例。

\paragraph{效果。}
数组元素的正确读写是精确的——
在索引合法（$0 \leq i < N$）的前提下，
标量化和 BTOR2 array 的读写值与 C 语义一致。
小数组标量化额外提供了零初始化的保证。

\paragraph{局限性。}
（1）\textbf{越界访问不可检测}。
BTOR2 的 read/write 是全函数，
索引超出 $[0, N)$ 范围时仍返回一个值。
CBMC 前端注入的数组边界检查断言
（\texttt{property\_class == "array bounds"}）
被编码器\emph{跳过}，不生成 bad 属性。
这是已知局限——当前模型不验证数组越界安全性。

（2）\textbf{BTOR2 array 的初始值无约束}。
如果程序逻辑依赖全局/静态数组的零初始化，
且该数组未被标量化，
模型检查器可能因为初始值不确定而产生虚假反例。


% =====================================================================
\subsection{指针与内存对象}
\label{sec:pointer}

\paragraph{类型与分析。}
C 指针是最复杂也最容易出错的类型。
指针的核心操作包括：
取地址（$\&x$）、解引用（$*p$）、
指针算术（$p + n$，按元素步进）、
指针比较（$p == q$，$p < q$）、
指针类型转换（\texttt{(int*)p}）、
动态内存分配（\texttt{malloc} / \texttt{free}）。
C 的指针语义隐含一个\emph{字节级内存模型}：
内存是一维字节序列，指针是字节地址，
指针算术按字节计算，强制类型转换重新解释底层字节。
指针还涉及\emph{别名}（多个指针指向同一对象）、
\emph{悬空指针}（free 后使用）和\emph{空指针}解引用等问题。

\paragraph{BTOR2 的能力。}
BTOR2 完全没有指针概念，也没有字节级内存模型。
它只有位向量和数组，
不知道"地址"是什么。

\paragraph{困难。}
（1）\textbf{地址的表示}。
C 指针是字节地址，但 BTOR2 没有地址空间。
需要设计一个地址表示方案。

（2）\textbf{解引用的目标确定}。
给定一个指针 $p$，解引用 $*p$ 需要知道 $p$ 指向\emph{哪个对象的哪个位置}。
但 BTOR2 数组是独立的，没有统一的地址空间把对象串起来。

（3）\textbf{指针算术}。
$p + 1$ 在 C 中按元素大小步进，
但在对象 ID 模型中，对象之间没有顺序关系。

（4）\textbf{动态分配}。
\texttt{malloc} 在运行时创建新对象，
但 BTOR2 的状态变量在编码时就必须固定。

（5）\textbf{别名与不确定性}。
指针可能在运行时指向多个不同对象
（如 \texttt{p = c ? \&a : \&b}），
解引用时需要考虑所有可能的目标。

\paragraph{方法：对象标识模型 + 目标解析表。}

我们采用\emph{对象级}（而非字节级）的内存模型。
核心设计：

\textbf{指针表示。}
指针编码为位向量 $\lceil \texttt{T*} \rceil = \mathrm{bv}\langle w \rangle$（$w \geq 32$），
值是\emph{对象标识符}——一个分配给每个可寻址对象的唯一正整数常量。
对象 ID 的分配：
0 = NULL，1/2 保留给 CBMC 内部对象
（\texttt{\_\_CPROVER\_deallocated}、\texttt{\_\_CPROVER\_dead\_object}），
3+ 按需分配给用户对象。

\textbf{目标解析表。}
编码器维护一张映射表 $\mathrm{OT}$，
记录每个指针符号当前指向的后备对象和元素索引：
\[
\mathrm{OT}(p) = (\mathit{array\_id},\; \mathit{idx},\; \tau_{\mathrm{elem}},\; \tau_{\mathrm{storage}},\; \mathit{stride})
\]
其中 $\mathit{array\_id}$ 是后备 BTOR2 数组（或标量化数组）的符号标识，
$\mathit{idx}$ 是元素索引表达式，
$\mathit{stride}$ 是展平后的元素步长（用于指针算术）。

\textbf{取地址} $\&x$：
查 $x$ 的后备对象，返回常量 $\mathit{const}(\mathrm{id}(x))$。
若 $x$ 是数组元素 $a[i]$，
则记录 $\mathrm{OT}$ 中索引为 $i$，返回 $\mathit{const}(\mathrm{id}(a))$。

\textbf{解引用读} $*p$：
\begin{enumerate}
\item 查 $\mathrm{OT}(p)$ 得到 $(\mathit{array\_id},\; \mathit{idx})$。
\item 若后备是标量化数组，构建 ITE 选择器。
\item 若后备是 BTOR2 array state，生成 $\mathit{read}(\mathit{state}[\mathit{array\_id}],\; \mathit{idx})$。
\end{enumerate}

\textbf{解引用写} $*p := e$：
同理翻译为 $\mathit{write}$ 或逐元素条件更新。

\textbf{条件指针} $p = c\,?\,\&a : \&b$：
$\mathrm{OT}$ 记录多个 $(\mathit{guard},\; \mathit{target})$ 对。
$*p$ 翻译为嵌套 ITE：
$\ite(c,\; \mathit{read}(s_a, i_a),\; \mathit{read}(s_b, i_b))$。

\textbf{指针算术} $p + n$：
查 $\mathrm{OT}(p)$ 得到 $(\mathit{array\_id},\; \mathit{idx},\; \mathit{stride})$，
新索引 $= \mathit{idx} + n \times \mathit{stride}$，
对象不变。
这只在同一数组内部有效——
跨对象的指针算术在对象 ID 模型中没有意义。

\textbf{malloc} 处理：
为每次 \texttt{malloc} 调用创建合成数组
$\_\mathtt{heap\_malloc\$}k : \mathit{state}\;\mathit{array}\langle 64,\; w_{\mathrm{elem}} \rangle$，
注册到 $\mathrm{OT}$，
返回指针值 $= \mathit{const}(\mathrm{id}(\_\mathtt{heap\_malloc\$}k))$。
若分配类型是结构体，按字段投影创建多个后备数组。
循环内的 \texttt{malloc} 被检测并标记为歧义指针，
强制上近似处理。
\texttt{free} 不产生状态效应，仅推进 PC。

\textbf{字符串常量}：
字符串字面量 \texttt{"hello"} 获得合成后备数组
$\_\mathtt{string\_literal\$hello}$，
读取时按字符位置构建 ITE 选择器。

\textbf{解析失败的处理}：
当指针无法静态解析时
（如来自函数参数、函数指针调用的返回值、
或指向两个以上对象的歧义指针）：
读取回退为无约束 input（$\mathit{fresh\_unknown}$）；
写入尝试广播到所有已知目标（保守上近似），
若无已知目标则静默丢弃。

\paragraph{效果。}
对\emph{可静态解析}的指针——
指向编译时已知的全局/局部对象、索引是常量或简单表达式——
读写翻译是精确的。
条件指针（有限个可能目标）通过 ITE 链精确处理。
\texttt{malloc} 的堆对象被正确建模为独立数组。

\paragraph{局限性。}
指针编码是三层中保真度最低的，存在以下已知限制：

（1）\textbf{不是字节级模型}。
不维护内存布局，不支持 \texttt{memcpy}、
\texttt{char*} 逐字节读写、
通过指针的类型双关（type-punning via pointer）。
元素步长 $\mathit{stride}$ 是展平后的叶子计数，不是字节数。

（2）\textbf{指针算术受限}。
同一数组内的指针算术可以处理，
但跨对象的指针减法（$p - q$）、
任意地址运算不支持。

（3）\textbf{别名分析保守}。
当指针指向 2 个以上对象时被标记为歧义，
解引用回退为上近似。
不做数据流分析来精化别名关系。

（4）\textbf{悬空指针与 use-after-free}。
\texttt{free} 不产生效应，堆对象状态永久存在。
CBMC 的 \texttt{\_\_CPROVER\_deallocated} 追踪机制部分支持
（保留 ID 1/2），但不保证完整。

（5）\textbf{解析失败时的可靠性}。
指针写入被静默丢弃时，
程序后续通过该指针观察到的不变量可能被遗漏，
可靠性不保证。
这是编码器在精度和可终止性之间做的权衡——
遇到无法精确处理的指针时选择继续编码而非终止报错。

\subsection{形式语义与可靠性}
\label{sec:formal-semantics}

本节给出编码的形式化基础。
我们首先定义 Goto-program 的小步操作语义，
然后定义 BTOR2 迁移系统的语义，
建立两者之间的状态对应关系，
最后证明可靠性定理：
若 BTOR2 系统判定 safe，则原始程序在受支持片段上满足所有断言。

% =====================================================================
\subsubsection{源语义：Goto-program 小步操作语义}
\label{sec:source-sem}

\paragraph{配置。}
Goto-program 的运行时状态是一个二元组 $(\ell,\; \sigma)$，其中：
\begin{itemize}
\item $\ell \in \{0, 1, \ldots, n\}$ 是\emph{位置编号}，
对应第~\ref{sec:pc}~节中分配的 PC 值；
\item $\sigma : \mathit{Var} \to \mathit{Val}$ 是\emph{存储}，
将每个（展平后的）变量标识符映射到其当前值。
对于标量化数组，$\sigma$ 中的键是 $a\$0, a\$1, \ldots$；
对于结构体字段，键是 $s.f_1, s.f_2, \ldots$；
对于 BTOR2 数组状态，$\sigma(a)$ 是一个从索引到元素值的函数。
\end{itemize}

\paragraph{表达式求值。}
给定存储 $\sigma$，表达式 $e$ 的值记为 $\llbracket e \rrbracket_\sigma$，
按 C 语义递归定义：
\begin{align*}
\llbracket c \rrbracket_\sigma &= v(c)
  && \textit{(常量)} \\
\llbracket v \rrbracket_\sigma &= \sigma(v)
  && \textit{(变量读取)} \\
\llbracket e_1 \oplus e_2 \rrbracket_\sigma
  &= \llbracket e_1 \rrbracket_\sigma \;\hat\oplus\; \llbracket e_2 \rrbracket_\sigma
  && \textit{($\hat\oplus$ 为 C 语义下的运算)} \\
\llbracket a[i] \rrbracket_\sigma &= \sigma(a)(\llbracket i \rrbracket_\sigma)
  && \textit{(数组索引)} \\
\llbracket s.f \rrbracket_\sigma &= \sigma(s.f)
  && \textit{(字段访问，展平后)} \\
\llbracket *p \rrbracket_\sigma &= \sigma(\mathrm{OT}(p).\mathit{array})(\mathrm{OT}(p).\mathit{idx})
  && \textit{(指针解引用，对象级模型)}
\end{align*}

\paragraph{迁移规则。}
小步迁移关系 $(\ell, \sigma) \to (\ell', \sigma')$ 由位置 $\ell$ 处的指令类型决定。
以下用推理规则给出，其中 $\mathrm{instr}(\ell)$ 表示位置 $\ell$ 处的 Goto 指令。

\medskip
\noindent\textbf{赋值。}
\[ \dfrac{
  \mathrm{instr}(\ell) = (x := e) \quad
  v = \llbracket e \rrbracket_\sigma
}{
  (\ell,\; \sigma) \to (\ell{+}1,\; \sigma[x \mapsto v])
} \;\;\textsc{(Assign)}
\]
赋值将 $x$ 在 $\sigma$ 中的值更新为右值 $e$ 的求值结果，
位置推进到 $\ell+1$。
对于复合左值（$a[i]$、$s.f$、$*p$），$x$ 是展平后的目标标识符，
更新语义相应投影。

\medskip
\noindent\textbf{无条件跳转。}
\[ \dfrac{
  \mathrm{instr}(\ell) = \texttt{goto}\;t
}{
  (\ell,\; \sigma) \to (t,\; \sigma)
} \;\;\textsc{(Goto)}
\]

\medskip
\noindent\textbf{条件跳转。}
\[
\dfrac{
  \mathrm{instr}(\ell) = \texttt{goto}(g,\; t) \quad
  \llbracket g \rrbracket_\sigma = \mathit{true}
}{
  (\ell,\; \sigma) \to (t,\; \sigma)
} \;\;\textsc{(Goto-T)}
\]
\[
\dfrac{
  \mathrm{instr}(\ell) = \texttt{goto}(g,\; t) \quad
  \llbracket g \rrbracket_\sigma = \mathit{false}
}{
  (\ell,\; \sigma) \to (\ell{+}1,\; \sigma)
} \;\;\textsc{(Goto-F)}
\]
条件跳转有两个后继：守卫 $g$ 为真时跳到 $t$，为假时顺序执行 $\ell+1$。
这涵盖了 C 的 \texttt{if}、\texttt{while}、\texttt{for}、\texttt{switch}
降低后的所有跳转模式。

\medskip
\noindent\textbf{假设。}
\[
\dfrac{
  \mathrm{instr}(\ell) = \texttt{assume}\;c \quad
  \llbracket c \rrbracket_\sigma = \mathit{true}
}{
  (\ell,\; \sigma) \to (\ell{+}1,\; \sigma)
} \;\;\textsc{(Assume)}
\]
当 $\llbracket c \rrbracket_\sigma = \mathit{false}$ 时，\textsc{Assume} 规则不适用——
配置\emph{卡住}（stuck），该路径终止，不再产生后续状态。
这对应 C 语义中 \texttt{assume} 排除不可行路径的效果。

\medskip
\noindent\textbf{断言。}
\[
\dfrac{
  \mathrm{instr}(\ell) = \texttt{assert}\;c
}{
  (\ell,\; \sigma) \to (\ell{+}1,\; \sigma)
} \;\;\textsc{(Assert)}
\]
断言\emph{不阻塞}执行——无论条件是否成立，配置都推进到 $\ell+1$。
断言违反通过\emph{坏状态}集合独立定义（见下文）。

\paragraph{坏状态与安全性。}
源程序的坏状态集合定义为：
\[
\mathrm{Bad}_G = \{(\ell,\; \sigma) \mid
  \mathrm{instr}(\ell) = \texttt{assert}(c) \;\wedge\;
  \llbracket c \rrbracket_\sigma = \mathit{false} \}
\]
初始配置为 $(0,\; \sigma_0)$，其中 $\sigma_0$ 由变量的静态初始值或零初始化决定。

\begin{definition}[源程序安全性]
\label{def:source-safety}
程序 $G$ 是\emph{安全的}，当且仅当不存在有限迁移序列
$(0, \sigma_0) \to^* (\ell, \sigma)$ 使得 $(\ell, \sigma) \in \mathrm{Bad}_G$。
\end{definition}

% =====================================================================
\subsubsection{目标语义：BTOR2 迁移系统}
\label{sec:target-sem}

\paragraph{配置。}
BTOR2 迁移系统的运行时状态是 $(\mathit{pc},\; s)$，其中：
\begin{itemize}
\item $\mathit{pc} \in \{0, 1, \ldots, n\}$ 是程序计数器的当前值
（编码为 $\lceil \log_2(n{+}1) \rceil$ 位位向量）；
\item $s : X \to \mathit{BV}$ 是对所有 BTOR2 状态变量的赋值，
将每个状态变量映射到其位向量值。
\end{itemize}

\paragraph{迁移。}
BTOR2 的迁移是完全确定的——每一步的所有状态更新同时发生：
\[
(\mathit{pc},\; s) \to (\mathit{pc}',\; s')
\]
其中：
\begin{align*}
\mathit{pc}' &= \llbracket \mathit{next}(\mathit{pc}) \rrbracket_s \\
s'(x) &= \llbracket \mathit{next}(x) \rrbracket_s \quad \text{对每个状态变量 } x \in X \setminus \{\mathit{pc}\}
\end{align*}
所有 $\mathit{next}$ 函数都在\emph{当前}状态 $s$ 上求值（同步语义）。
$\mathit{next}$ 函数是嵌套 ITE 表达式，其求值结果由 $\mathit{pc}$ 的当前值
（决定哪个守卫 $\mathit{at}(\ell)$ 为真）和变量值共同决定。

\paragraph{约束与坏状态。}
BTOR2 系统的约束和坏状态：
\begin{align*}
\mathit{feasible}(\mathit{pc}, s) &\iff
  \bigwedge_i \llbracket \mathit{constraint}_i \rrbracket_{(\mathit{pc}, s)} = \mathit{true} \\
\mathrm{Bad}_M &= \{(\mathit{pc}, s) \mid
  \bigvee_{\ell \in \mathit{Asserts}} (\mathit{pc} = \ell \;\wedge\;
  \llbracket \neg c_\ell \rrbracket_s) \}
\end{align*}
其中 $\mathit{Asserts}$ 是所有 \texttt{assert} 位置的集合，
$c_\ell$ 是位置 $\ell$ 处的断言条件。

\begin{definition}[BTOR2 系统安全性]
\label{def:target-safety}
BTOR2 系统 $M$ 是\emph{安全的}，当且仅当不存在有限迁移序列
$(0, s_0) \to^* (\mathit{pc}, s)$ 使得 $(\mathit{pc}, s) \in \mathrm{Bad}_M$，
且路径上所有中间配置都满足 $\mathit{feasible}$。
\end{definition}

% =====================================================================
\subsubsection{状态对应}
\label{sec:correspondence}

我们定义源配置与目标配置之间的\emph{对应关系} $\sim$，
它是可靠性证明的核心桥梁。

首先定义值编码函数 $\mathrm{enc}_\tau : \mathit{Val}_\tau \to \mathit{BV}$，
将 C 类型 $\tau$ 的值映射为位向量：
对标量类型，$\mathrm{enc}$ 是数值的二进制表示
（有符号数用补码，浮点用 IEEE 754 打包位模式）。

\begin{definition}[状态对应]
\label{def:correspondence}
$(\ell,\; \sigma) \sim (\mathit{pc},\; s)$ 当且仅当：
\begin{enumerate}
\item \textbf{位置一致}：$\mathit{pc} = \ell$。

\item \textbf{标量变量一致}：
  对每个标量状态变量 $x_v$，
  $s(x_v) = \mathrm{enc}_{\tau_v}(\sigma(v))$。

\item \textbf{结构体字段一致}：
  对每个展平后的字段状态 $x_{v.f}$，
  $s(x_{v.f}) = \mathrm{enc}_{\tau_f}(\sigma(v.f))$。

\item \textbf{标量化数组一致}：
  对每个标量化数组元素 $x_{a\$i}$，
  $s(x_{a\$i}) = \mathrm{enc}(\sigma(a)(i))$。

\item \textbf{BTOR2 数组一致}：
  对每个 BTOR2 数组状态 $a$，
  $\forall j \in [0, N).\;
  s(a)(j) = \mathrm{enc}(\sigma(a)(j))$。
\end{enumerate}
\end{definition}

注意条件 5 只要求合法索引范围内一致——
越界索引处的值在源语义中是未定义的，
在目标语义中是 BTOR2 全函数的任意返回值，
两者不可比较，也不需要比较。

% =====================================================================
\subsubsection{表达式对应}
\label{sec:expr-correspondence}

\begin{lemma}[表达式对应]
\label{lem:expr}
设 $(\ell, \sigma) \sim (\mathit{pc}, s)$。
对任意受支持类型上的 C 表达式 $e$，
\[
\llbracket e \rrbracket^{M}_s
= \mathrm{enc}_{\tau_e}\!\left(\llbracket e \rrbracket_\sigma\right)
\]
其中 $\llbracket \cdot \rrbracket^M_s$ 是 BTOR2 表达式在状态 $s$ 上的求值。
\end{lemma}

\noindent\textit{证明思路。}
对表达式结构做归纳：

\begin{itemize}
\item \textbf{基础情况}：
  常量 $c$ 的 BTOR2 编码 $\mathit{const}(\mathrm{enc}(v(c)))$
  在任何状态下求值都是 $\mathrm{enc}(v(c))$。
  变量 $v$ 的读取直接由定义~\ref{def:correspondence} 的条件 2--5 保证。

\item \textbf{算术运算}（$+$, $-$, $\times$）：
  BTOR2 的 $\mathit{add/sub/mul}$ 在位向量上的行为
  与 C 的回绕语义一致（模 $2^n$ 运算）。

\item \textbf{除法与取模}：
  符号性在操作选择时已正确处理
  （$\mathit{sdiv}$ 对有符号，$\mathit{udiv}$ 对无符号）。
  当除数非零时，BTOR2 的除法与 C 语义一致。

\item \textbf{比较}：
  有符号比较使用 $\mathit{slt/sgt/sgte}$，
  无符号使用 $\mathit{ult/ugt/ugte}$，
  与 C 的比较语义一致。

\item \textbf{类型转换}：
  宽化时 $\mathit{sext/uext}$ 的选择与源类型符号性一致；
  窄化时 $\mathit{slice}$ 取低位，与 C 截断语义一致。

\item \textbf{复合左值读取}（$a[i]$、$s.f$、$*p$）：
  标量化数组的 ITE 选择器、BTOR2 数组的 $\mathit{read}$、
  指针的目标解析，
  分别由定义~\ref{def:correspondence} 的条件 3--5 保证其读取的底层值一致，
  再由索引 / 字段表达式的归纳假设完成证明。
\end{itemize}

% =====================================================================
\subsubsection{迁移对应与可靠性}
\label{sec:simulation}

\begin{lemma}[单步对应]
\label{lem:step}
设 $(\ell, \sigma) \sim (\mathit{pc}, s)$，且 $\mathrm{instr}(\ell)$ 是受支持的指令类型。
若 $(\ell, \sigma) \to (\ell', \sigma')$，
则存在 $s'$ 使得 $(\mathit{pc}, s) \to (\mathit{pc}', s')$，
且 $(\ell', \sigma') \sim (\mathit{pc}', s')$。
\end{lemma}

\noindent\textit{证明。}
对指令类型做分情况讨论。

\smallskip
\noindent\textbf{情况 ASSIGN}（$x := e$）：
源语义给出 $(\ell, \sigma) \to (\ell{+}1, \sigma[x \mapsto v])$，
其中 $v = \llbracket e \rrbracket_\sigma$。
在目标系统中，$\mathit{state\_updates}[x]$ 包含对 $(\mathit{at}(\ell), \llbracket e \rrbracket^M)$，
$\mathit{pc\_updates}$ 包含 $(\mathit{at}(\ell), \ell{+}1)$。
由于当前 $\mathit{pc} = \ell$，守卫 $\mathit{at}(\ell)$ 为真，
$\mathit{next}$ 函数求值为：
\begin{align*}
s'(x) &= \llbracket e \rrbracket^M_s
  = \mathrm{enc}(\llbracket e \rrbracket_\sigma)
  && \text{(引理~\ref{lem:expr})}
  = \mathrm{enc}(v) = \mathrm{enc}(\sigma'(x)) \\
\mathit{pc}' &= \ell + 1
\end{align*}
其余未更新的状态变量保持不变（$\mathit{next}$ 的默认值是当前值），
因此 $s'$ 与 $\sigma'$ 满足定义~\ref{def:correspondence}。

\smallskip
\noindent\textbf{情况 GOTO}（\texttt{goto $t$}）：
$\mathit{pc\_updates}$ 包含 $(\mathit{at}(\ell), t)$，
故 $\mathit{pc}' = t$。存储不变，对应关系保持。

\smallskip
\noindent\textbf{情况 GOTO-T / GOTO-F}：
$\mathit{pc\_updates}$ 同时包含
$(\mathit{at}(\ell) \wedge g, t)$ 和
$(\mathit{at}(\ell) \wedge \neg g, \ell{+}1)$。
由引理~\ref{lem:expr}，$g$ 在 $s$ 上的求值与在 $\sigma$ 上一致，
因此正确的分支被选中。

\smallskip
\noindent\textbf{情况 ASSUME}：
当 $\llbracket c \rrbracket_\sigma = \mathit{true}$ 时，
$\mathit{pc\_updates}$ 中的 $(\mathit{at}(\ell) \wedge c, \ell{+}1)$ 生效，
PC 推进。
同时，$\mathit{constraint}_\ell = \neg\mathit{at}(\ell) \vee c$ 在此配置上为真，
不违反可行性。
当 $\llbracket c \rrbracket_\sigma = \mathit{false}$ 时，
源配置卡住（无迁移），
引理的前提不满足，无需证明。

\smallskip
\noindent\textbf{情况 ASSERT}：
$\mathit{pc\_updates}$ 包含 $(\mathit{at}(\ell), \ell{+}1)$，
PC 无条件推进。存储不变，对应关系保持。 \hfill $\square$

\begin{theorem}[可靠性]
\label{thm:soundness}
设程序 $G$ 仅使用受支持的指令和类型构造，
且 $M$ 是其 BTOR2 编码。
若 $M$ 是安全的（定义~\ref{def:target-safety}），
则 $G$ 是安全的（定义~\ref{def:source-safety}）。
\end{theorem}

\noindent\textit{证明。}
反证法。假设 $G$ 不安全，
则存在迁移序列
$(0, \sigma_0) \to^* (\ell, \sigma)$
使得 $(\ell, \sigma) \in \mathrm{Bad}_G$，
即 $\mathrm{instr}(\ell) = \texttt{assert}(c)$
且 $\llbracket c \rrbracket_\sigma = \mathit{false}$。

对初始配置 $(0, \sigma_0)$ 和 $M$ 的初始配置 $(0, s_0)$，
由 init 的构造（零初始化和静态初值的翻译），
$(0, \sigma_0) \sim (0, s_0)$ 成立。

对迁移序列的每一步应用引理~\ref{lem:step}，
得到 $M$ 中的对应序列
$(0, s_0) \to^* (\ell, s)$，
且 $(\ell, \sigma) \sim (\ell, s)$。
路径上所有中间配置满足 $\mathit{feasible}$
（每个 assume 的 constraint 在对应配置上成立，
由 ASSUME 情况的论证保证）。

由定义~\ref{def:correspondence} 和引理~\ref{lem:expr}，
$\llbracket c \rrbracket^M_s = \mathrm{enc}(\llbracket c \rrbracket_\sigma) = \mathrm{enc}(\mathit{false})$，
故 $\llbracket \neg c \rrbracket^M_s = \mathit{true}$，
因此 $(\ell, s) \in \mathrm{Bad}_M$。

这与 $M$ 安全矛盾。 \hfill $\square$

\begin{remark}[虚假反例]
\label{rem:spurious}
定理~\ref{thm:soundness} 的逆\emph{不}成立：
$M$ 不安全\emph{不}蕴含 $G$ 不安全。
上近似引入了源程序中不存在的状态
（如 BTOR2 数组状态的无约束初始值、
无法解析的指针读取的无约束返回值、
不支持表达式的无约束替换），
这些状态可能触发 bad 属性，产生\emph{虚假反例}（spurious counterexample）。
第~\ref{sec:experiments}~节将量化虚假反例的比例。
\end{remark}

\begin{remark}[不受支持构造的影响]
\label{rem:unsupported}
定理~\ref{thm:soundness} 的前提是程序仅使用受支持的构造。
当程序包含不受支持的构造时，
可靠性保证可能在以下情况被削弱：
\begin{itemize}
\item 变量类型不受支持时，该变量被完全跳过（无状态、无更新），
  引理~\ref{lem:step} 的对应关系不成立。
\item 指针写入被静默丢弃时，
  后续通过该指针观察到的状态可能与源程序不一致。
\item 残余内存操作（\texttt{havoc\_object} 等）被忽略时，
  可能遗漏对程序状态的副作用。
\end{itemize}
在实践中，编码器对不受支持的构造发出编译警告，
实验中统计其出现频率（第~\ref{sec:experiments}~节）。
\end{remark}


\subsection{标记迁移系统的定义}

转换器将 goto 程序编码为五元组 $M = (\State, \Init, \Trans, \Bad, \Cons)$：

\begin{definition}[BTOR2 迁移系统]
  \hfill
  \begin{itemize}
    \item \textbf{状态空间} $\State = \Loc \times \Val_{v_1} \times \cdots \times \Val_{v_k}$，
          其中 $\Loc = \{0, 1, \ldots, N{-}1\}$ 为程序位置集合，
          $v_1, \ldots, v_k$ 为收集到的程序变量。
          状态 $\sigma = (\ell, x_1, \ldots, x_k) \in \State$，
          其中 $\ell$ 为程序计数器值，$x_i$ 为变量 $v_i$ 的当前值。
    \item \textbf{初始状态集} $\Init \subseteq \State$。
    \item \textbf{迁移关系} $\Trans \subseteq \State \times \Input \times \State$，
          允许外部输入 $\iota \in \Input$（对应 nondet 调用）。
    \item \textbf{坏状态集} $\Bad \subseteq \State$，编码断言违反。
    \item \textbf{约束} $\Cons : \State \times \Input \to \mathbb{B}$，
          编码 \texttt{assume} 语句。
  \end{itemize}
\end{definition}

\subsection{位置谓词}

核心构件是\textbf{位置谓词}，表示"程序计数器当前在位置 $\ell$"：
\begin{equation}
  \guard_\ell(\sigma) \;\triangleq\; [\,\sigma.\mathit{pc} = \ell\,]
\end{equation}

转换器把"在位置 $\ell$ 执行指令 $i$" 编码为"该步的更新被 $\guard_\ell$ 守卫"。
对应代码 \texttt{make\_pc\_equals(loc)}（\texttt{goto\_btor2.cpp:227}）。

\subsection{初始状态}

\begin{definition}[初始状态集]
  \[
    \Init = \{\, (0,\;\mathit{init}_1,\;\ldots,\;\mathit{init}_k) \,\}
  \]
  其中程序计数器初始值为 $0$（程序入口），变量初始值：
  \[
    \mathit{init}_i =
    \begin{cases}
      0 & \text{标量变量（} \texttt{builder.zero} \text{）} \\
      \bot \text{（任意值）} & \text{非标量化数组状态}
    \end{cases}
  \]
\end{definition}

\begin{remark}
  非标量化数组（$>256$ 元素的数组）不生成 \texttt{init} 指令，在 $t{=}0$
  取任意值。对依赖隐式零初始化的全局数组，这会丢失约束，可能导致伪反例。
\end{remark}

\subsection{迁移关系：逆序 ITE 链}

每个变量 $v$ 维护一个\textbf{更新集} $U_v = [(c_1, e_1),\ldots,(c_m, e_m)]$
（按源程序顺序）。PC 亦然。下一状态由逆序 ITE 链给出
（\texttt{goto\_btor2.cpp:258--297}）：

\begin{definition}[单步迁移]
  \label{def:transition}
  对变量 $v$，其下一状态为：
  \[
    \nextv{v}(\sigma, \iota) =
    \begin{cases}
      e_j(\sigma, \iota) & \text{若 } j = \max\{\,i \mid c_i(\sigma,\iota) = \True\,\} \text{ 存在} \\
      v(\sigma) & \text{若所有 } c_i \text{ 均为假（保持当前值）}
    \end{cases}
  \]
  PC 的下一状态 $\nextpc$ 同理。
  一步迁移 $(\sigma, \iota) \to \sigma'$ 定义为对所有 $v$：
  $\sigma'.v = \nextv{v}(\sigma, \iota)$ 且 $\sigma'.\mathit{pc} = \nextpc(\sigma, \iota)$。
\end{definition}

\begin{lemma}[守卫互斥性]
  \label{lem:mutex}
  对同一变量 $v$，$U_v$ 中任意两个守卫 $c_i, c_j$（$i \neq j$）满足
  $c_i \wedge c_j \equiv \False$。
\end{lemma}

\begin{proof}
  每个守卫 $c_i$ 形如 $\guard_\ell$ 或 $\guard_\ell \wedge g$。不同更新的
  $\guard_\ell$ 来自不同位置 $\ell$（每条指令只 push 一次更新，且
  \texttt{process\_*} 提前 \texttt{return} 保证单路径）。由于 $\mathit{pc}$
  取唯一值，$\guard_{\ell_1} \wedge \guard_{\ell_2} \equiv \False$（$\ell_1 \neq \ell_2$）。
  同一 $\ell$ 内的多个更新仅出现在：
  \begin{itemize}
    \item 条件 GOTO：$g$ 与 $\neg g$，互斥。
    \item 标量化数组 store：不同 \texttt{is\_match}$_i$（索引匹配），互斥。
    \item 歧义指针写入：写入所有目标，但各目标的更新集 entry 独立，互斥性由 ITE 链保持。
  \end{itemize}
\end{proof}

\begin{corollary}
  ITE 链的构造顺序对 $\nextv{v}$ 的语义无影响。代码中"后者覆盖前者"的
  逆序遍历是防御性冗余，不改变语义。
\end{corollary}

\subsection{可达集}

\begin{definition}[约束下的可达集]
  \[
    \Reach = \mu\, X.\; (\Init \cap \Cons) \;\cup\;
    \{\,\sigma' \mid \exists \sigma \in X,\; \iota.\;\;
       \Trans(\sigma, \iota, \sigma') \wedge \Cons(\sigma', \iota)\,\}
  \]
  其中 $\mu$ 表示最小不动点。
\end{definition}

\subsection{各类指令的语义方程}

表~\ref{tab:semantics} 给出每种指令对 $(\Trans, \Bad, \Cons)$ 的增量贡献。
$\eval{e}$ 表示表达式 $e$ 在当前状态下的求值（BTOR2 组合电路）。

\begin{table}[h]
  \centering
  \small
  \caption{各类指令的语义方程}
  \label{tab:semantics}
  \renewcommand{\arraystretch}{1.4}
  \begin{tabular}{>{\ttfamily}l >{$}l<{$} l}
    \toprule
    \textrm{指令} & \textrm{语义增量} & \textrm{代码位置} \\
    \midrule
    x := e & U_v \mathrel{+}= \{(\guard_\ell,\; \eval{e})\} & assign.cpp \\
    goto T (uncond) & U_{pc} \mathrel{+}= \{(\guard_\ell,\; T)\} & goto \\
    if g goto T & U_{pc} \mathrel{+}= \{(\guard_\ell\wedge g,\; T),\;(\guard_\ell\wedge\neg g,\; \ell{+}1)\} & goto \\
    assert(c) & \Bad \mathrel{+}= \{\sigma \mid \guard_\ell(\sigma) \wedge \neg\eval{c}(\sigma)\} & assert \\
    assume(c) & \Cons \mathrel{\wedge\!=} (\neg\guard_\ell \vee \eval{c}) & assume \\
    x = nondet() & U_v \mathrel{+}= \{(\guard_\ell,\; \iota_{\mathit{fresh}})\} & func\_call \\
    abort() & \Cons \mathrel{\wedge\!=} \neg\guard_\ell & func\_call \\
    \bottomrule
  \end{tabular}
\end{table}

% =========================================================================
\section{可靠性定理}
\label{sec:soundness}
% =========================================================================

\subsection{可靠性定理陈述}

\begin{theorem}[转换可靠性]
  \label{thm:sound}
  设 $P$ 为满足前提 \textbf{P1}--\textbf{P7} 的 C 程序（见下），$\varphi$
  为编码为 $\Bad$ 集的安全性质。若 BTOR2 模型 $M$ 满足 $\varphi$
  （即 $\Reach_M \cap \Bad = \emptyset$），则 C 程序 $P$ 也满足 $\varphi$：
  \[
    M \models \varphi \;\Longrightarrow\; P \models \varphi
  \]
\end{theorem}

\textbf{前提条件}：
\begin{description}[font=\bfseries,labelwidth=2em]
  \item[P1] 程序已被 goto 层内联为单一 \texttt{main}，无过程调用残留。
  \item[P2] 所有 GOTO 指令至多一个 target（无 switch jump table），或转换器正确处理了多目标。
  \item[P3] 所有指针解引用的目标可被静态解析为单一后备数组，或为\textbf{歧义指针}（已被可靠过近似）。
  \item[P4] 循环内的 \texttt{malloc} 不依赖"每次分配返回独立内存"的语义，或被降级为 nondet。
  \item[P5] 非标量化数组（$>256$ 元素）不依赖隐式零初始化。
  \item[P6] 不使用 \texttt{START\_THREAD}、\texttt{THROW}、\texttt{CATCH} 等未建模指令。
  \item[P7] 所有未支持的指令类型被降级为保守的过近似（nondet 或 fallthrough），而非静默丢弃语义。
\end{description}

\subsection{证明}

\begin{proof}[定理~\ref{thm:sound} 的证明（草图）]
  分四步。

  \medskip\noindent\textbf{步骤 1：单步保持（赋值公理）。}
  对每条指令 $i$ 在位置 $\ell$，由引理~\ref{lem:mutex}（守卫互斥性），
  其效果 $\nextv{v}$ 在 $\guard_\ell$ 时恰好执行 $i$ 的语义。
  对于赋值 $x := e$，$\nextv{x} = \eval{e}$，与霍尔赋值公理
  $\{P[x \mapsto e]\}\;x{:=}e\;\{P\}$ 一致。

  \medskip\noindent\textbf{步骤 2：组合保持（霍尔复合规则）。}
  由 $\guard_\ell$ 互斥，一步只执行一条指令，故多步执行 = 指令序列的霍尔复合：
  \[
    \frac{\{P\}\;S_1\;\{Q\} \quad \{Q\}\;S_2\;\{R\}}
         {\{P\}\;S_1;S_2\;\{R\}}
  \]
  每条指令贡献一条 $\guard_\ell$ 守卫的更新，互不干扰。

  \medskip\noindent\textbf{步骤 3：约束等价（assume 的双重编码）。}
  \texttt{assume(c)} 产生双重效应（\texttt{goto\_btor2.cpp:341--360}、
  \texttt{goto\_btor2\_instructions.cpp:448--468}）：
  \begin{itemize}
    \item 约束 $\Cons \wedge= (\neg\guard_\ell \vee \eval{c})$ 把
          $\mathit{pc}=\ell \wedge \neg c$ 的状态从可达集剔除。
    \item PC 守卫 $\guard_\ell \wedge \eval{c}$ 保证从 $\ell$ 仅在 $c$ 真时转移到 $\ell{+}1$。
  \end{itemize}
  两者合起来等价于霍尔语义 $\{c \wedge P\}\;S\;\{Q\}$ 的路径过滤。

  \medskip\noindent\textbf{步骤 4：过近似不破坏可靠性。}
  \begin{itemize}
    \item \texttt{abort} 排除 $\mathit{pc}=\ell$ 的状态（过近似）。
    \item 未解析指针读返回 fresh nondet input（扩大可能值域）。
    \item 歧义指针写写入所有目标（覆盖真实写入位置的超集）。
    \item \texttt{malloc} 摘要为单一堆对象（扩大别名可能）。
  \end{itemize}
  这些操作\textbf{只扩大}可达集或\textbf{只减少}路径，不产生"虚假证明"
  （即不会把不满足 $\varphi$ 的程序误判为满足）。因此
  $M \models \varphi \Rightarrow P \models \varphi$ 成立。
\end{proof}

\begin{remark}
  反方向（完备性）\textbf{不成立}：BTOR2 报违反时，反例可能是过近似导致的
  \emph{假反例}（spurious counterexample），需 witness 验证。
\end{remark}

\subsection{循环与归纳}

循环在 BTOR2 中编码为迁移系统中的环（回边 $\ell_{\mathrm{tail}} \to \ell_{\mathrm{head}}$）。
霍尔逻辑的 While 规则需要人工提供不变式：
\[
  \frac{\{I \wedge b\}\;S\;\{I\}}
       {\{I\}\;\texttt{while}\;b\;\texttt{do}\;S\;\{I \wedge \neg b\}}
\]

BTOR2 编码\textbf{不计算} $I$，而是把整个循环交给模型检测器：
\begin{itemize}
  \item \textbf{BMC}：展开到深度 $k$，找 $\leq k$ 步的反例。
  \item \textbf{$k$-induction}：自动发现归纳不变量。
  \item \textbf{IC3/PDR}：基于子句的归纳推理。
\end{itemize}

转换器的职责是\textbf{忠实编码迁移关系} $\Trans$，循环不变量的发现
\textbf{完全委托给下游}归纳引擎。理论上 $\Trans$ 允许任意长轨迹；
实际能否证明取决于模型检测器的归纳能力。

% =========================================================================
\section{语义层次：操作语义与霍尔逻辑的关系}
\label{sec:semantic-level}
% =========================================================================

\subsection{操作语义 vs.\ 指称语义}

\begin{definition}[操作语义（Operational Semantics）]
  操作语义通过\textbf{状态迁移规则}定义程序的含义：
  $\langle S, \sigma \rangle \to \langle S', \sigma' \rangle$。
  程序的意义就是它执行的步骤序列。
\end{definition}

\begin{definition}[指称语义（Denotational Semantics）]
  指称语义通过\textbf{数学对象}（函数、域元素）定义程序的含义：
  $\llbracket S \rrbracket : \State \rightharpoonup \State$。
  程序的意义是一个从输入状态到输出状态（或结果）的部分函数。
\end{definition}

\begin{remark}
  转换器编码的是\textbf{操作语义}，具体来说是\textbf{小步结构化操作语义}
  （Structural Operational Semantics, SOS）。

  \textbf{关键证据}：BTOR2 的迁移关系 $\Trans$ 对应单步归约规则
  $\langle S, \sigma \rangle \to \langle S', \sigma' \rangle$。
  每条 goto 指令 = 一步迁移；程序计数器 $\mathit{pc}$ 显式编码"下一条要执行
  的指令"。这是典型的 small-step 配置迁移
  （configuration transition），其中配置是 $(\mathit{pc}, \sigma)$。

  \textbf{为什么不是大步语义}：大步语义（Natural Semantics / Big-Step）
  的规则形如 $\langle S, \sigma \rangle \Downarrow \sigma'$，一条规则
  消化整个语句（如整个循环）。BTOR2 模型\textbf{没有}一步执行整个循环的
  迁移；循环被编码为多步迁移形成的环，由模型检测器展开。因此是\textbf{小步}。
\end{remark}

\subsection{霍尔逻辑与操作语义的关系}

霍尔逻辑（Hoare Logic）是\textbf{公理语义}（Axiomatic Semantics），
与操作语义处于不同的抽象层次：

\begin{center}
\renewcommand{\arraystretch}{1.3}
\begin{tabular}{>{\bfseries}l l l}
  \toprule
  语义层次 & 研究对象 & 在本项目中的角色 \\
  \midrule
  操作语义 & 单步迁移 $\to$ & 转换器\textbf{生成}的 BTOR2 模型 \\
  公理语义 & 三元组 $\{P\}S\{Q\}$ & 用于\textbf{证明}模型的正确性 \\
  \bottomrule
\end{tabular}
\end{center}

两者\textbf{互补}：
\begin{itemize}
  \item BTOR2 是\textbf{模型}（model）：描述"程序做什么"。
  \item 霍尔逻辑是\textbf{证明系统}（proof system）：描述"如何推导程序的性质"。
\end{itemize}

\subsection{可以写出霍尔推理规则吗？}

\textbf{可以}。虽然 BTOR2 编码的是操作语义，但我们可以从迁移关系 $\Trans$
导出一组霍尔推理规则，用于\textbf{验证}转换的正确性。表~\ref{tab:hoare}
给出从 BTOR2 编码导出的霍尔规则（对应 \S\ref{sec:formal-semantics} 的语义方程）。

\begin{table}[h]
  \centering
  \small
  \caption{从 BTOR2 编码导出的霍尔推理规则}
  \label{tab:hoare}
  \renewcommand{\arraystretch}{1.8}
  \begin{tabular}{l}
    \toprule
    \textbf{赋值公理}（对应 \texttt{state\_updates} + \texttt{at\_loc}）： \\
    $\{P[x \mapsto \eval{e}]\}\;x {:=} e\;\{P\}$ \\
    \midrule
    \textbf{顺序复合}（对应 fall-through PC 更新）： \\
    $\dfrac{\{P\}\;S_1\;\{Q\} \quad \{Q\}\;S_2\;\{R\}}
           {\{P\}\;S_1;\,S_2\;\{R\}}$ \\
    \midrule
    \textbf{条件分支}（对应 \texttt{process\_goto} 的双守卫）： \\
    $\dfrac{\{g \wedge P\}\;S_T\;\{Q\} \quad \{\neg g \wedge P\}\;S_{\ell+1}\;\{Q\}}
           {\{P\}\;\texttt{if}\;g\;\texttt{goto}\;T\;\{Q\}}$ \\
    \midrule
    \textbf{循环}（对应回边；归纳委托给模型检测器）： \\
    $\dfrac{\{I \wedge g\}\;S_{\mathrm{body}}\;\{I\}}
           {\{I\}\;\texttt{while}\;g\;\texttt{do}\;S\;\{I \wedge \neg g\}}$ \\
    \midrule
    \textbf{断言}（对应 \texttt{bad} 谓词）： \\
    $\dfrac{}{\{P\}\;\texttt{assert}(c)\;\{P \wedge c\}}$
       \quad（若 $\Bad$ 不触发） \\
    \midrule
    \textbf{假设}（对应 \texttt{constraint} + PC 守卫）： \\
    $\dfrac{}{\{P\}\;\texttt{assume}(c)\;\{P \wedge c\}}$ \\
    \bottomrule
  \end{tabular}
\end{table}

\begin{remark}
  这些霍尔规则的\textbf{可靠性}（soundness）恰好由定理~\ref{thm:sound}
  保证：若 BTOR2 模型满足某性质，则对应的霍尔三元组在该程序中成立。
  反之不成立（不完备），因为过近似可能引入假反例。
\end{remark}

\subsection{小步操作语义的推导规则}
\label{sec:sos-rules}

转换器生成的 BTOR2 模型本质上是一组小步归约规则。下面给出其完整的
结构化操作语义（SOS）推导规则。配置（configuration）为
$C = (\ell, \sigma, \iota)$，其中 $\ell$ 为 PC 值，$\sigma$ 为变量状态，
$\iota$ 为当前步的外部输入。归约关系 $C \to C'$ 表示一步迁移。

\subsubsection*{基本记号}

\begin{itemize}
  \item $\sigma[x \mapsto v]$：状态 $\sigma$ 中变量 $x$ 更新为值 $v$。
  \item $\eval{e}_{\sigma,\iota}$：表达式 $e$ 在状态 $\sigma$、输入 $\iota$ 下的求值。
  \item $\ell{+}1$：PC 顺序后继。
  \item $\mathit{fresh}(\tau)$：类型 $\tau$ 的新鲜 nondet 输入值。
\end{itemize}

\subsubsection*{推导规则}

\begin{mathpar}
  \inferrule*[left=\textsc{Assign}]
    { }
    {(\ell,\;\sigma,\;\iota) \;\to\; (\ell{+}1,\;\sigma[x \mapsto \eval{e}_{\sigma,\iota}],\;\iota')}
  \quad\text{（当 $\ell$ 处指令为 $x {:=} e$）}

  \inferrule*[left=\textsc{Goto-True}]
    { }
    {(\ell,\;\sigma,\;\iota) \;\to\; (T,\;\sigma,\;\iota')}
  \quad\text{（当 $\ell$ 处为 \texttt{goto T}，无条件跳转）}

  \inferrule*[left=\textsc{Goto-Taken}]
    { \eval{g}_{\sigma,\iota} = \True }
    {(\ell,\;\sigma,\;\iota) \;\to\; (T,\;\sigma,\;\iota')}
  \quad\text{（当 $\ell$ 处为 \texttt{if g goto T}，条件为真）}

  \inferrule*[left=\textsc{Goto-Fall}]
    { \eval{g}_{\sigma,\iota} = \False }
    {(\ell,\;\sigma,\;\iota) \;\to\; (\ell{+}1,\;\sigma,\;\iota')}
  \quad\text{（当 $\ell$ 处为 \texttt{if g goto T}，条件为假）}

  \inferrule*[left=\textsc{Assert-Ok}]
    { \eval{c}_{\sigma,\iota} = \True }
    {(\ell,\;\sigma,\;\iota) \;\to\; (\ell{+}1,\;\sigma,\;\iota')}
  \quad\text{（断言成功：PC 推进）}

  \inferrule*[left=\textsc{Assert-Fail}]
    { \eval{c}_{\sigma,\iota} = \False }
    {(\ell,\;\sigma,\;\iota) \;\to\; (\ell,\;\sigma,\;\iota')}
  \quad\text{（断言失败：PC 停滞，$\Bad$ 触发）}

  \inferrule*[left=\textsc{Assume-Ok}]
    { \eval{c}_{\sigma,\iota} = \True }
    {(\ell,\;\sigma,\;\iota) \;\to\; (\ell{+}1,\;\sigma,\;\iota')}
  \quad\text{（假设成立：PC 推进）}

  \inferrule*[left=\textsc{Assume-Fail}]
    { \eval{c}_{\sigma,\iota} = \False }
    {(\ell,\;\sigma,\;\iota) \;\not\to}
  \quad\text{（假设不成立：无后继迁移，路径终止）}

  \inferrule*[left=\textsc{Nondet}]
    { }
    {(\ell,\;\sigma,\;\iota) \;\to\; (\ell{+}1,\;\sigma[x \mapsto \mathit{fresh}(\tau)],\;\iota')}
  \quad\text{（$x = \mathit{nondet}()$：$x$ 获得任意值）}

  \inferrule*[left=\textsc{Abort}]
    { }
    {(\ell,\;\sigma,\;\iota) \;\to\; (\ell,\;\sigma,\;\iota')}
  \quad\text{（abort：终态自环，PC 永久停滞）}

  \inferrule*[left=\textsc{Fallthrough}]
    { }
    {(\ell,\;\sigma,\;\iota) \;\to\; (\ell{+}1,\;\sigma,\;\iota')}
  \quad\text{（DECL/SKIP/LOCATION 等：仅 PC 递增）}
\end{mathpar}

\subsubsection*{多步执行与轨迹}

一条执行轨迹（trace）是配置序列 $C_0 \to C_1 \to C_2 \to \cdots$，其中
$C_0 \in \Init$。性质 $\varphi$ 成立当且仅当所有轨迹都不进入 $\Bad$ 状态：

\[
  P \models \varphi \;\iff\; \forall \text{ trace } C_0 \to C_1 \to \cdots
  .\; \forall k.\; C_k \notin \Bad
\]

模型检测器的任务是对这条轨迹的\textbf{所有}有限前缀和（理想地）无限延伸
进行推理，通常通过 BMC（有界展开）或 IC3（归纳推理）完成。

\begin{remark}
  这些 SOS 规则与表~\ref{tab:hoare} 的霍尔规则之间存在直接的对应：
  每条霍尔规则都可以从对应的 SOS 规则\textbf{推导}出来。例如，
  \textsc{Assert-Ok} 和 \textsc{Assert-Fail} 两条 SOS 规则蕴含
  霍尔规则 $\{P\}\;\texttt{assert}(c)\;\{P \wedge c\}$
  （因为失败路径停滞，成功路径保证 $c$）。这种对应关系是
  \emph{操作语义与公理语义等价}的经典实例。
\end{remark}

% =========================================================================
\section{内存模型的支持程度}
\label{sec:memory-model}
% =========================================================================

\subsection{现有内存模型的刻画}

BTOR2 无原生指针类型。转换器采用\textbf{对象级（object-based）内存模型}，
而非字节级（byte-addressable）内存模型。内存被表示为：

\begin{definition}[对象级内存模型]
  内存由一组命名的\textbf{后备数组}（backing arrays）组成：
  \[
    \mathit{Mem} = \{\, \mathit{obj}_i : \mathit{Array}(\mathit{Idx}_i, \mathit{Elem}_i) \,\}
  \]
  每个可寻址对象（全局变量、局部数组、\texttt{malloc} 堆对象）对应一个
  后备数组。指针值 $p$ 编码为 $(\mathit{obj\_id}, \mathit{offset})$，
  其中 $\mathit{obj\_id}$ 是对象的唯一整数标识。
\end{definition}

\begin{center}
\textbf{表：内存操作的支持程度}
\end{center}
\renewcommand{\arraystretch}{1.3}
\begin{tabular}{>{\bfseries}l l l}
  \toprule
  内存操作 & 支持程度 & 说明 \\
  \midrule
  取地址 \&x & 完全支持 & 编码为 $\mathit{obj\_id}$ 常量 \\
  解引用 *p（单目标指针） & 完全支持 & 解析为 BTOR2 \texttt{read} \\
  解引用 *p（歧义指针） & 可靠过近似 & 读取返回 nondet \\
  通过指针写（单目标） & 完全支持 & 解析为 BTOR2 \texttt{write} \\
  通过指针写（歧义） & 可靠过近似 & 写入\textbf{所有}已知目标 \\
  指针算术 p + n & 部分支持 & 仅索引偏移；对象切换不建模 \\
  数组越界访问 & 不检查 & \texttt{array bounds} 断言被跳过 \\
  指针比较 & 支持 & 基于 $\mathit{obj\_id}$ 比较 \\
  malloc 分配 & 摘要支持 & 单一堆对象；循环内不独立 \\
  空指针解引用 & 不建模 & 无 NULL 检查 \\
  字节级访问（byte\_extract） & 有限支持 & 偏移为 0 时正确 \\
  指针类型转换 & 支持 & element\_type 跟随 cast 目标 \\
  \bottomrule
\end{tabular}

\subsection{不完备支持的正确表述}

\textbf{结论}：转换器\textbf{没有显式的、完备的内存模型}，但提供了
\textbf{不完备但可靠（sound）的内存抽象}。

\begin{shaded*}
\textbf{准确表述}：
转换器实现了\textbf{对象级别名分析}（object-level alias analysis），
而非完整的字节级内存模型。对于指针指向可静态解析的程序，内存操作
被精确建模；对于指针指向不可静态解析的程序（如歧义指针、循环内
重赋值），内存操作被降级为可靠的过近似（nondet 读取 + 全目标写入）。
\end{shaded*}

\textbf{缺失的内存模型组件}：
\begin{enumerate}
  \item \textbf{无空间安全检查}：数组越界、指针越界不被检测
        （\texttt{array bounds} 断言被显式跳过）。
  \item \textbf{无 NULL 指针检查}：空指针解引用不产生错误状态。
  \item \textbf{无 use-after-free}：\texttt{free} 后的访问不被建模。
  \item \textbf{无字节级别名}：不同类型指针通过同一地址访问不同大小的数据
        （\texttt{union} 的核心语义）不精确建模。
  \item \textbf{无完整堆模型}：\texttt{malloc} 返回单一摘要对象，
        循环内多次分配不独立。
\end{enumerate}

这些缺失意味着：对\textbf{内存安全性质}（buffer overflow、null deref、
use-after-free），当前编码\textbf{不完备}；但对\textbf{功能正确性性质}
（如算法是否正确实现），在 P3--P5 成立时编码可靠。

% =========================================================================
\section{已知限制与修复状态}
\label{sec:limitations}
% =========================================================================

\begin{center}
\textbf{表：已识别问题及修复状态}
\end{center}
\renewcommand{\arraystretch}{1.3}
\begin{tabular}{>{\bfseries}c p{4cm} p{3.5cm} p{3cm}}
  \toprule
  \# & 问题描述 & 严重性 & 状态 \\
  \midrule
  1 & 流不敏感指针映射（循环内重赋值导致假阴性） & 严重—不可靠 & 已修复：歧义指针检测 + 全目标写入 \\
  2 & 循环内 malloc 摘要为单一对象 & 中等—精度损失 & 待修复 \\
  3 & 多目标 GOTO 只取首目标 & 中等—语义丢失 & 待修复 \\
  4 & abort 编码为约束而非路径截断 & 低—安全属性可靠 & 待修复 \\
  5 & assert 失败后控制流继续 & 低—安全属性可靠 & 待修复 \\
  \bottomrule
\end{tabular}

% =========================================================================
\section{总结}
% =========================================================================

\begin{enumerate}
  \item 转换器编码的是\textbf{小步操作语义}（SOS），每条 goto 指令对应
        一步迁移，PC 显式编码控制流位置。
  \item 该编码可导出一组\textbf{霍尔推理规则}（表~\ref{tab:hoare}），
        其可靠性由定理~\ref{thm:sound} 保证。
  \item 内存模型为\textbf{对象级别名分析}，不是完备的字节级内存模型，
        但在指针目标可静态解析时精确，否则可靠过近似。
  \item 循环的推理\textbf{委托给下游}模型检测器的归纳引擎
        （BMC / k-induction / IC3）。
\end{enumerate}

\end{document}
